\section{Maximum likelihood estimation}
\label{sec:estimation}

With the transition density in hand, estimation is the maximization of
\eqref{eq:loglik} over
$\thetavec = (\mu, \sigma, p, \theta_J)$. The estimator class,
\code{JumpDiffusionEstimator}, wraps this in a numerically defensive
objective and offers two optimizers whose contrast is one of the library's
central themes: a fast local method that needs a good start, and a global
method that needs none.

\subsection{The objective}

Optimizers minimize, so the objective is the \emph{negative}
log-likelihood, with two guard layers wrapped around \eqref{eq:loglik}:

\begin{lstlisting}[caption={The objective in \code{estimation/maximum\_likelihood.py}.}]
def log_likelihood(self, params: np.ndarray) -> float:
    sigma, jump_prob = params[1], params[2]

    if sigma <= 0:
        return np.inf
    if jump_prob < 0 or jump_prob > 1:
        return np.inf

    self._model.update_parameters(**dict(zip(self._param_names, params)))
    value = self._model.log_likelihood(self.increments, self.dt)
    # Penalize non-numeric likelihood values so that population-based
    # optimizers simply discard the offending candidates -- same device
    # as in Ospina Arango (2009).
    if not np.isfinite(value):
        return np.inf
    return -value
\end{lstlisting}

\begin{decision}[Map every pathology to $+\infty$]
\label{dec:penalize}
Invalid parameters ($\sigma \le 0$, $p \notin [0,1]$) and non-finite
likelihood evaluations (numerical accidents deep in the FFT or in extreme
tails) all return $+\infty$. For a population-based optimizer this is
exactly the right poison: an infinite objective makes the candidate lose
every selection comparison and disappear from the population without
corrupting the ranking of the finite candidates. The same device ---
mapping \code{NaN} evaluations to a large positive constant --- was used
in \citet{Ospina2009}; $+\infty$ is its cleanest float representation.
\end{decision}

\subsection{Starting values from moments}
\label{sec:initial-guess}

Gradient methods need a starting point. The estimator builds one from the
sample moments of the increments: since
$\E[\Delta X] = \mu \Delta t + p\,\E[J]$ and
$\operatorname{Var}[\Delta X] = \sigma^2 \Delta t + p \operatorname{Var}[J]
+ p(1-p)\E[J]^2$, ignoring the (typically small) jump contributions gives
the first-pass guesses
\begin{equation}
  \hat\mu_0 = \frac{\overline{\Delta x}}{\Delta t},
  \qquad
  \hat\sigma_0 = \frac{s_{\Delta x}}{\sqrt{\Delta t}},
  \qquad
  \hat p_0 = 0.1,
  \label{eq:initial-guess}
\end{equation}
and each jump family contributes its own heuristic through its
\code{initial\_guess} method --- e.g.\ the skew-normal
starts \code{jump\_scale} at half the increment standard deviation and
signs \code{jump\_skew} by the sample skewness. These are deliberately
crude: their job is to land in the right order of magnitude, not to be
accurate. When they are not enough --- and Section~\ref{sec:de} shows this
happens in practice --- the answer is a global optimizer, not a cleverer
guess.

\subsection{Local optimization: L-BFGS-B}

The default method hands the objective to
\code{scipy.optimize.minimize} with \code{method="L-BFGS-B"}, box
constraints from the model ($\sigma > 0$, $0 < p < 1$, plus each jump
parameter's bounds from Section~\ref{sec:contract}), and tolerances
\code{maxiter=1000}, \code{ftol=1e-9}. On well-behaved likelihoods (normal
or skew-normal jumps, moderate $p$) it converges in milliseconds. Its
failure modes on this problem are well documented in
\citet{Ospina2009}: with heavy-tailed or strongly skewed jump laws the
mixture likelihood develops multiple local optima and long, curved ridges,
and a quasi-Newton iterate either stalls on a box boundary or polishes
whichever local optimum captures it.

\subsection{Global optimization: differential evolution}
\label{sec:de}

Differential evolution \citep{StornPrice1997} maintains a population of
$\mathrm{NP}$ candidate vectors and evolves it by mutation, crossover and
greedy selection. One generation of the \emph{rand/1/bin} strategy --- the
one used here --- transforms each population member $x_i$ as follows:
\begin{align}
  \text{mutation:}\quad
    & v_i = x_{r_1} + F\,(x_{r_2} - x_{r_3}),
    \qquad r_1 \ne r_2 \ne r_3 \ne i \ \text{random},
    \label{eq:de-mutation}\\[2pt]
  \text{crossover:}\quad
    & u_{i,d} =
      \begin{cases}
        v_{i,d} & \text{if } U_d \le \mathrm{CR} \ \text{or}\ d = d_{\mathrm{rand}},\\
        x_{i,d} & \text{otherwise},
      \end{cases}
    \label{eq:de-crossover}\\[2pt]
  \text{selection:}\quad
    & x_i \leftarrow u_i \ \text{iff}\ f(u_i) \le f(x_i).
    \label{eq:de-selection}
\end{align}
Mutation \eqref{eq:de-mutation} perturbs a random member by a scaled
difference of two others, so step sizes adapt to the population's own
spread in each direction --- large early on, shrinking as the population
contracts onto a basin. Binomial crossover \eqref{eq:de-crossover} mixes
mutant and parent coordinate-wise. Selection \eqref{eq:de-selection} is
greedy but elitist: the best-ever point can never be lost. No gradients
appear anywhere, so multimodality, ridges, and the occasional
$+\infty$ from Decision~\ref{dec:penalize} are all handled gracefully.

\begin{decision}[Port the thesis' \code{DEoptim} configuration]
\label{dec:de-config}
The applied finding of \citet{Ospina2009} --- refined in
\citet{Ardia2011} --- is that DE with R's \code{DEoptim} defaults
(strategy rand/1, $\mathrm{NP} = 70$, 400 generations) recovers the SGED
jump-diffusion optimum even from very wide bounds, where L-BFGS-B and
simulated annealing both fail. The library ports this configuration to
\code{scipy.optimize.differential\_evolution}:
\begin{itemize}
  \item \code{strategy="rand1bin"} is exactly rand/1 with binomial
        crossover;
  \item SciPy sizes the population as $\mathrm{popsize} \times k$ where
        $k$ is the number of parameters, so \code{popsize=10} reproduces
        $\mathrm{NP} = 70$ for the 7-parameter SGED model and scales
        proportionally for other jump families;
  \item \code{maxiter=400} matches the generation cap, and SciPy's
        convergence tolerance (\code{tol=0.01}, stop when the population's
        objective spread falls below 1\% of its mean) implements the
        early-stopping criterion the thesis recommends --- in practice
        runs stop well before 400 generations.
\end{itemize}
\end{decision}

\begin{decision}[Finite, data-driven search box]
\label{dec:finite-bounds}
DE samples its initial population uniformly over the bounds, so the
half-open boxes used by L-BFGS-B ($\mu$ unbounded, $\sigma < \infty$)
must be replaced by finite ones. The library builds them from the data:
$\mu$ ranges over $\hat\mu_0 \pm \max(1,\ 10|\hat\mu_0|,\ 10\hat\sigma_0)$,
$\sigma$ over $(10^{-6},\ 10\,\hat\sigma_0)$, and any unbounded jump
parameter is capped at $\pm 20\, s_{\Delta x}$ --- a single jump twenty
times the typical increment is a generous ceiling, especially since
$s_{\Delta x}$ is itself inflated by whatever jumps the data contains.
Wide boxes are deliberate: the thesis' point is precisely that DE does
not need tight prior knowledge, so the defaults encode ignorance, not
information. Bounds the user supplies, and finite bounds from the jump
family, are kept as-is.
\end{decision}

\begin{lstlisting}[caption={DE configuration in \code{estimation/maximum\_likelihood.py} (excerpt).}]
if bounds is None:
    bounds = self._finite_bounds()
lows, highs = zip(*bounds)
de_bounds = Bounds(np.asarray(lows, float), np.asarray(highs, float))

de_kwargs = {
    "strategy": "rand1bin",   # DEoptim's rand/1 with binomial crossover
    "maxiter": 400,           # thesis generation cap
    "popsize": 10,            # 10 * 7 params = NP 70 for the SGED model
}
de_kwargs.update(kwargs)      # user overrides, e.g. seed=, workers=
return differential_evolution(func=self.log_likelihood,
                              bounds=de_bounds, **de_kwargs)
\end{lstlisting}

The cost asymmetry is worth stating plainly: L-BFGS-B typically needs
tens-to-hundreds of likelihood evaluations, DE thousands --- seconds
instead of milliseconds with closed-form convolutions, and proportionally
more through the FFT path. The recommended workflow is therefore to fit
with DE (optionally passing \code{seed=} for reproducibility and
\code{workers=-1} to parallelize the population evaluation) and reserve
L-BFGS-B for quick exploratory fits and for the repeated re-fits inside
bootstrap loops, where the data are simulated near a known optimum.

\subsection{Reported quantities}

\code{estimate()} returns the fitted parameters, the maximized
log-likelihood $\loglik(\hat\thetavec)$, and the information criteria
\begin{equation}
  \mathrm{AIC} = 2k - 2\loglik(\hat\thetavec),
  \qquad
  \mathrm{BIC} = k \log n - 2\loglik(\hat\thetavec),
  \label{eq:aic-bic}
\end{equation}
with $k = \dim\thetavec$ and $n$ the number of increments --- the
ingredients for the model comparison of Section~\ref{sec:comparison} ---
plus the raw optimizer diagnostics and a convergence flag that downstream
bootstrap loops use to discard failed re-fits.
